\chapter{课堂练习及结果}
\label{cp:introduction}
% TODO: 在这里写第一章内容，替换掉示例文字。
% 这是一行测试的文字\\换行后如何？Welcome to the introduction of your dissertation. The introduction of a dissertation serves as a critical component, setting the tone and laying the foundation for the entire research endeavour. It is tasked with providing a clear and concise overview of the research topic, elucidating the context and significance of the study within the broader academic landscape. A well-crafted dissertation introduction should delineate the research problem or question, offering a rationale for its relevance and addressing any existing gaps in knowledge. Furthermore, it typically outlines the objectives and aims of the study, guiding the reader through the anticipated contributions and outcomes. In addition, the introduction often encapsulates the methodology employed, presenting the chosen approach and rationale behind it. Lastly, it functions as a road-map, offering a brief glimpse into the structure and organisation of the dissertation, thereby orienting the reader and facilitating comprehension of the subsequent chapters. Overall, a dissertation introduction should engage the reader's interest, provide a clear framework for the research, and justify its importance in the academic realm. For a clearer and reader-friendly experience on referencing chapters, please refer to the chapter titled \nameref{cp:citations} (referred to as \autoref{cp:citations}).
\section{第 1 题 含空格文件名的批量整理}
\subsection{实验内容}

首先，如图\ref{fig:make_dirs}所示，创建 input/docs 和 input/tmp；

\begin{figure}[H]
    \centering
    \includegraphics[width=0.8\textwidth]{截屏2026-08-29 19.14.49.png}
    \caption{创建 input/docs 和 input/tmp}
    \label{fig:make_dirs}
\end{figure}

如图\ref{fig:创建文件}所示，使用vim编辑器创建 input/docs/notes one.txt（内容为 alpha 和 beta 两行）、
input/docs/.secret.txt（内容为 hidden）、input/tmp/empty.txt（空文件）

\begin{figure}[H]
    \centering
    \includegraphics[width=0.8\textwidth]{截屏2026-08-29 19.15.21.png}
    \caption{创建 input/docs/notes one.txt、input/docs/.secret.txt、input/tmp/empty.txt}
    \label{fig:创建文件}
\end{figure}

如图\ref{fig:执行脚本}所示，先使用vim编辑器创建空的run.log，然后使用date > run.log和date >> run.log命令将当前日期分两行写入 run.log，最后使用cat run.log命令查看 run.log 的内容。

\begin{figure}[H]
    \centering
    \includegraphics[width=0.8\textwidth]{截屏2026-08-29 19.16.12.png}
    \caption{执行脚本}
    \label{fig:执行脚本}
\end{figure}

如图\ref{fig:绝对路径}所示，使用pwd路径显示q01的绝对路径。

\begin{figure}[H]
    \centering
    \includegraphics[width=0.8\textwidth]{截屏2026-08-29 19.16.46.png}
    \caption{绝对路径}
    \label{fig:绝对路径}
\end{figure}

如图\ref{fig:列出项目}所示，先用find命令找出 input 目录下的所有项目，然后对于每个找到的项目使用ls命令列出其详细信息。

\begin{figure}[H]
    \centering
    \includegraphics[width=0.8\textwidth]{截屏2026-08-29 19.17.34.png}
    \caption{列出项目}
    \label{fig:列出项目}
\end{figure}

在图\ref{fig:列出项目}中，find为查找命令，“."表示查找当前目录，-type f表示查找的为文件，-exec表示对每个找到的项目执行后面的命令，{}表示当前找到的项目，\;表示-exec命令的结束。
而在exec后的ls为列出命令，-l为长格式选项，-a为显示所有文件选项，会显示隐藏文件。

如图\ref{fig:复制文件}所示，利用管道将find和while结合起来，递归遍历当前目录下所有`.txt`文件，在 \verb|work/<25020007076>| 目录里完整复刻原有的文件夹结构，把全部 txt 文件复制过去。

\begin{figure}[H]
    \centering
    \includegraphics[width=0.8\textwidth]{截屏2026-08-29 22.21.54.png}
    \caption{复制文件}
    \label{fig:复制文件}
\end{figure}

在图\ref{fig:复制文件}中，IFS表示关闭字符串分割，防止文件名包含空格、制表符被拆分，read -r file表示读取每一行的文件名，while read -r file; do ... done表示对每个读取到的文件名执行do和done之间的命令，mkdir -p表示创建目录，dirname \$file表示获取文件所在的目录，cp \$file \verb|work/<25020007076>/$file| 表示将文件复制到目标目录。

如图\ref{fig:设置权限}所示，将复制后的目录权限设为 750、普通文件权限设为 640；生成 inventory.txt，列出相对路径和字节数。

\begin{figure}[H]
    \centering
    \includegraphics[width=0.8\textwidth]{截屏2026-08-29 22.28.11.png}
    \caption{设置权限}
    \label{fig:设置权限}
\end{figure}

在图\ref{fig:设置权限}中，find为查找命令，“."表示查找当前目录，-type d表示查找的为目录，-exec表示对每个找到的项目执行后面的命令，{}表示当前找到的项目，\;表示-exec命令的结束。而在exec后的chmod为修改权限命令，750表示目录权限为rwxr-x---，640表示文件权限为rw-r-----。最后使用find命令查找所有文件并使用ls -l命令列出详细信息，并将结果输出到 inventory.txt 文件中。

如图\ref{fig:查看结果}所示，使用cat命令查看 inventory.txt 文件的内容。

\begin{figure}[H]
    \centering
    \includegraphics[width=0.8\textwidth]{截屏2026-08-29 22.28.25.png}
    \caption{查看结果}
    \label{fig:查看结果}
\end{figure}

\subsection{实验结果}

第一题的github的commit记录和push记录分别如图\ref{fig:github记录1}所示。

\begin{figure}[H]
    \centering
    \includegraphics[width=0.8\textwidth]{截屏2026-08-29 22.31.42.png}
    \caption{github记录1}
    \label{fig:github记录1}
\end{figure}

\section{第 2 题 随机访问日志统计}

\subsection{实验内容}

如图\ref{fig:输入csv记录}所示，使用vim编辑器创建access.csv文件，并输入如图\ref{fig:csv内容}所示的内容。

\begin{figure}[H]
    \centering
    \includegraphics[width=0.8\textwidth]{截屏2026-08-30 11.00.39.png}
    \caption{输入csv记录}
    \label{fig:输入csv记录}
\end{figure}

\begin{figure}[H]
    \centering
    \includegraphics[width=0.8\textwidth]{截屏2026-08-30 10.44.20.png}
    \caption{csv内容}
    \label{fig:csv内容}
\end{figure}

如图\ref{fig:analyze.sh}所示，使用vim编辑器编写 analyze.sh，并用cat命令查看 analyze.sh 

\begin{figure}[H]
    \centering
    \includegraphics[width=0.8\textwidth]{截屏2026-08-30 11.00.53.png}
    \caption{analyze.sh}
    \label{fig:analyze.sh}
\end{figure}

在图\ref{fig:analyze.sh}中，先判断传入的地址是否存在，如果不存在则将错误写入 stderr 并返回非零退出码。其中 [ ! -f \verb|"$1"| ] 为判断条件。\verb|>&2| 将错误消息输出到标准错误 stderr。
在输入“HTTP 5xx 数量最多的前 2 个 path:”的提示后，使用 awk 解析 CSV，跳过表头，统计状态码 5xx 的各 path 出错次数，按计数降序排序后取前 2 条输出。
在输入“平均 \texttt{latency\_ms:} ”的提示后，awk 跳过 CSV 表头，累加延迟字段并统计行数，计算得到保留两位小数的平均延迟并打印。

\begin{figure}[H]
    \centering
    \includegraphics[width=0.8\textwidth]{截屏2026-08-30 11.01.13.png}
    \caption{使用chmod命令为analyze.sh添加执行权限}
    \label{fig:添加analyze.sh权限}
\end{figure}   

如图\ref{fig:添加analyze.sh权限}所示，使用chmod命令为analyze.sh添加执行权限。然后如图\ref{fig:执行analyze.sh}所示，分别对 access.csv 和一个不存在的文件运行脚本analyze.sh，并查看输出结果。

\begin{figure}[H]
    \centering
    \includegraphics[width=0.8\textwidth]{截屏2026-08-30 11.01.27.png}
    \caption{执行analyze.sh}
    \label{fig:执行analyze.sh}
\end{figure}  

\subsection{实验结果}

第二题的github的commit记录和push记录分别如图\ref{fig:github记录2}与图\ref{fig:github记录3}所示。

\begin{figure}[H]
    \centering
    \includegraphics[width=0.8\textwidth]{截屏2026-08-30 11.01.54.png}
    \caption{github记录2}
    \label{fig:github记录2}
\end{figure}

\begin{figure}[H]
    \centering
    \includegraphics[width=0.8\textwidth]{截屏2026-08-30 11.02.04.png}
    \caption{github记录3}
    \label{fig:github记录3}
\end{figure}

\section{第 3 题 制造、解决并解释一次合并冲突}

\subsection{实验内容}

如图\ref{fig:初始化仓库}所示，使用git init命令初始化一个新的Git仓库。

\begin{figure}[H]
    \centering
    \includegraphics[width=0.8\textwidth]{截屏2026-08-31 08.27.04.png}
    \caption{初始化仓库}
    \label{fig:初始化仓库}
\end{figure}

如图\ref{fig:创建config}所示，创建config文件并输入内容mode=normal。

\begin{figure}[H]
    \centering
    \includegraphics[width=0.8\textwidth]{截屏2026-08-31 08.36.09.png}
    \caption{创建config}
    \label{fig:创建config}
\end{figure}

如图\ref{fig:提交config}所示，完成初次提交。

\begin{figure}[H]
    \centering
    \includegraphics[width=0.8\textwidth]{截屏2026-08-31 08.36.09.png}
    \caption{提交config文件}
    \label{fig:提交config文件}
\end{figure}

如图\ref{fig:创建feature-a分支}所示，创建feature-a分支。

\begin{figure}[H]
    \centering
    \includegraphics[width=0.8\textwidth]{截屏2026-08-31 08.38.07.png}
    \caption{创建feature-a分支}
    \label{fig:创建feature-a分支}
\end{figure}

如图\ref{fig:提交feature-a分支}所示，将config文件的内容内容改为 mode=safe 并提交。

\begin{figure}[H]
    \centering
    \includegraphics[width=0.8\textwidth]{截屏2026-08-31 08.38.07.png}
    \caption{提交feature-a分支}
    \label{fig:提交feature-a分支}
\end{figure}

如图\ref{fig:创建feature-b分支}所示，先切换回main分支，在main分支上创建feature-b分支。

\begin{figure}[H]
    \centering
    \includegraphics[width=0.8\textwidth]{截屏2026-08-31 08.40.32.png}
    \caption{创建feature-b分支}
    \label{fig:创建feature-b分支}
\end{figure}

如图\ref{fig:提交feature-b分支}所示，将config文件的内容改为 mode=fast 并提交。

\begin{figure}[H]
    \centering
    \includegraphics[width=0.8\textwidth]{截屏2026-08-31 08.42.08.png}
    \caption{提交feature-b分支}
    \label{fig:提交feature-b分支}
\end{figure}

如图\ref{fig:合并feature-a分支}所示，切换回main分支，尝试将feature-a分支合并到main分支。从图\ref{fig:合并feature-a分支}可以看出，没有出现合并错误，成功合并了feature-a分支。

\begin{figure}[H]
    \centering
    \includegraphics[width=0.8\textwidth]{截屏2026-08-31 08.43.25.png}
    \caption{合并feature-a分支}
    \label{fig:合并feature-a分支}
\end{figure}

从图\ref{fig:合并feature-b分支}可以看出，尝试将feature-b分支合并到main分支时，从git status的输出中可以看出出现了合并冲突。
\begin{figure}[H]
    \centering
    \includegraphics[width=0.8\textwidth]{截屏2026-08-31 08.43.53.png}
    \caption{合并feature-b分支}
    \label{fig:合并feature-b分支}
\end{figure}

\begin{figure}[H]
    \centering
    \includegraphics[width=0.8\textwidth]{截屏2026-08-31 08.44.15.png}
    \caption{查看合并冲突}
    \label{fig:查看合并冲突}
\end{figure}

从图\ref{fig:查看合并冲突}可以看出，config文件的内容出现了冲突，<<<<<<< HEAD表示当前分支的内容，=======表示分隔符，>>>>>>> feature-b表示要合并的分支的内容。

如图\ref{fig:解决合并冲突}所示，解决冲突时保留 mode=safe，并另加一行 note=reviewed。

\begin{figure}[H]
    \centering
    \includegraphics[width=0.8\textwidth]{截屏2026-08-31 08.44.43.png}
    \caption{解决合并冲突}
    \label{fig:解决合并冲突}
\end{figure}

\begin{figure}[H]
    \centering
    \includegraphics[width=0.8\textwidth]{截屏2026-08-31 08.45.11.png}
    \caption{提交解决冲突后的config文件}
    \label{fig:提交解决冲突后的config文件}
\end{figure}

从图\ref{fig:提交解决冲突后的config文件}可以看出，git status的输出显示已经解决了合并冲突，并且可以提交解决冲突后的config文件。

如图\ref{fig:查看提交记录}所示，确认工作区干净，运行 git log --all --graph --decorate --oneline 查看历史。

\begin{figure}[H]
    \centering
    \includegraphics[width=0.8\textwidth]{截屏2026-08-31 08.47.40.png}
    \caption{查看提交记录}
    \label{fig:查看提交记录}
\end{figure}

\subsection{实验结果}

如图\ref{fig:github记录4}所示，第三题的github commit记录和push记录。

\begin{figure}[H]
    \centering
    \includegraphics[width=0.8\textwidth]{截屏2026-08-31 09.24.35.png}
    \caption{github记录4}
    \label{fig:github记录4}
\end{figure}

\section{第 4 题 修复并构建一页技术说明}

\subsection{实验内容}

如图\ref{fig:创建report.tex文件}所示，使用vim编辑器创建report.tex文件，并补全report.tex文件的内容。再如图\ref{fig:查看report.tex文件}所示，使用cat命令查看report.tex文件的内容。
因使用题目要求的latexmk -pdf -halt-on-error无法编译中文，所以report.tex文件的内容使用英文表示。

\begin{figure}[H]
    \centering
    \includegraphics[width=0.8\textwidth]{截屏2026-08-30 16.35.21.png}
    \caption{创建report.tex文件}
    \label{fig:创建report.tex文件}
\end{figure}

\begin{figure}[H]
    \centering
    \includegraphics[width=0.8\textwidth]{截屏2026-08-31 09.29.39.png}
    \caption{查看report.tex文件}
    \label{fig:查看report.tex文件}
\end{figure}

如图\ref{fig:1使用latexmk编译report.tex文件}和图\ref{fig:2使用latexmk编译report.tex文件}所示，使用latexmk编译report.tex文件，并生成 report.pdf。

\begin{figure}[H]
    \centering
    \includegraphics[width=0.8\textwidth]{截屏2026-08-30 16.36.05.png}
    \caption{使用latexmk编译report.tex文件}
    \label{fig:1使用latexmk编译report.tex文件}
\end{figure}

\begin{figure}[H]
    \centering
    \includegraphics[width=0.8\textwidth]{截屏2026-08-30 16.37.26.png}
    \caption{使用latexmk编译report.tex文件}
    \label{fig:2使用latexmk编译report.tex文件}
\end{figure}

\subsection{实验结果}

如图\ref{fig:github记录5}所示，第四题的github commit记录和push记录。

\begin{figure}[H]
    \centering
    \includegraphics[width=0.8\textwidth]{截屏2026-08-30 16.37.53.png}
    \caption{github记录5}
    \label{fig:github记录5}
\end{figure}

\section{Github仓库地址}

本次课程的报告、练习、代码等，均已上传至个人github仓库。\\Github仓库地址：\url{https://github.com/lcx821/lcx_course_work}

其中，q03练习所用仓库也已上传至Github。\\仓库地址：\url{https://github.com/lcx821/q03}